\subsection{
    Матрица линейного оператора.
}

Пусть $\mathcal{V}$ и $\mathcal{W}$ - два линейных пространства.

Пусть $e = (\vec{e}_1, \ldots, \vec{e}_n)$ - некоторый базис в $\mathcal{V}$, $f = (\vec{f}_1, \ldots, \vec{f}_m)$ - некоторый базис в $\mathcal{W}$. 

Тогда $\mathscr{A}\vec{e}_1, \ldots, \mathscr{A}\vec{e}_n$ - это некоторые векторы в $\mathcal{W}$; значит, их можно, причем единственным образом, разложить по базису $\vec{f}_1, \ldots, \vec{f}_m$:

$$\mathscr{A}\vec{e}_1 = a_{11}\vec{f}_1 + \ldots + a_{m1}\vec{f}_m$$
$$\ldots \ldots \ldots \ldots \ldots \ldots \ldots \ldots \ldots$$
$$\mathscr{A}\vec{e}_n = a_{1n}\vec{f}_1 + \ldots + a_{mn}\vec{f}_m,$$

где $(a_{ij} \in \PP)$.

\begin{definition}
    Матрицу, составленную из координатных столбцов векторов $\mathscr{A}\vec{e}_1, \ldots, \mathscr{A}\vec{e}_n$ в базисе $f = (\vec{f}_1, \ldots, \vec{f}_m)$, называют \textit{\textbf{матрицей линейного оператора}} $\mathscr{A}$ в базисе $f$.
\end{definition}
